\section{Global Compatibility and the Inter-Face Sigma Frontier}

\subsection{Local law and global closure are distinct}

Project 41 also contains a stronger inter-face transport problem that
is logically separate from the native registered carrier used in the
relational-saturation theorem.

The seam construction supplies a locally admissible inter-face
transport vocabulary.

Within that construction, the local \(r_3\) action has a universal
form: it reverses side and polarity while preserving sheet.  Its
section routing contains both same-section and changed-section
branches.

The completed local vocabulary is also closed under the independently
derived neighbor-frame involution \(E\).

These are local and covariance results.

They do not by themselves construct a globally consistent seam
transport.

The remaining requirement for that stronger construction is that
selected seam actions satisfy the declared four-step closure relation
\[
(r_2r_3)^4=1
\]
throughout the complete flag system.

Thus the architecture continues to distinguish
\[
\text{local admissibility}
\]
from
\[
\text{global compatibility}.
\]

\subsection{The inter-face sigma}

\begin{definition}[Inter-face sigma]
An inter-face sigma is a globally selected assignment of locally
admissible seam transport relations across adjacent faces satisfying
all declared covariance and closure conditions.
\end{definition}

Sigma is not an additional long-range selector placed on top of the
local mechanics.

It is a compatibility relation for one particular stronger
inter-face transport construction.

Schematically,
\[
\text{local seam vocabulary}
\longrightarrow
\text{covariant compatibility}
\longrightarrow
\sigma
\longrightarrow
\text{admitted global inter-face transport}.
\]

\subsection{Completed E-closed seam vocabulary}

Neighbor-frame covariance enlarges the local admissible vocabulary.

After canonical completion under the involution \(E\), the global
search used
\[
2816
\]
distinct local seam options.

This completed vocabulary is the correct search domain for the
present global \(r_3\) problem.

The enlargement is mathematically important.

It shows that the earlier truncated vocabulary did not exhaust the
locally admissible seam possibilities once neighbor-frame covariance
was imposed.

But a larger locally admissible vocabulary does not imply global
closure.

\subsection{Completed Mac search}

The completed Mac solver run searched the \(2816\)-option
\(E\)-closed seam vocabulary for a globally consistent \(r_3\).

The run reached
\[
376
\]
total solver iterations and accumulated
\[
175069
\]
unique four-step closure nogoods.

No complete global \(r_3\) was found.

The search also did not prove the completed search space infeasible.

The result therefore remains logically open:
\[
\text{no solution found}
\]
and
\[
\text{no nonexistence theorem}.
\]

\subsection{Partial closure improvement}

The completed search nevertheless moved measurably deeper into the
closure constraints.

Earlier candidates typically failed the four-step relation on all
\[
1920
\]
flags.

Improved candidates first reached
\[
1904
\]
failed flags, corresponding to
\[
16
\]
correctly closing flags.

The completed Mac run reached repeated candidates with
\[
1888
\]
failed flags.

Thus the best observed partial closure increased to
\[
1920-1888=32
\]
correctly closing flags.

This does not constitute a global solution.

It does show that completion under neighbor-frame covariance admits
strictly deeper partial closure than the earlier truncated search.

\subsection{What the search establishes}

The Mac work supports the following internal conclusions:

\begin{enumerate}[label=(\roman*)]

\item local admissibility and global compatibility are distinct;

\item neighbor-frame covariance changes the admissible local search
space;

\item substantial local admissibility does not guarantee global
closure;

\item the unresolved global object is an inter-face transport
relation rather than an independent long-range selector;

\item the remaining obstruction concerns this stronger inter-face
construction, not the existence of the already certified registered
carrier.

\end{enumerate}

These conclusions constrain any future global transport extension.

They do not weaken the finite relational-saturation theorem proved on
the independently certified registered carrier.

\subsection{Sigma and a possible global summation domain}

Section 9 discusses lawful summation of finite contributions.

For a future inter-face construction, the sum cannot be taken over
arbitrary locally admissible seam terms.

Its admissible summation domain would have to be compatible with the
global inter-face transport relation.

Thus a possible future chain is
\[
\text{local seam contribution}
\longrightarrow
\sigma
\longrightarrow
\text{global inter-face summation domain}
\longrightarrow
\text{registered aggregate}.
\]

This remains a conditional extension.

It is not the mechanism used to prove native relational saturation in
Section 8.

\subsection{Circulation and winding}

If a globally compatible inter-face sigma is eventually constructed,
one may evaluate its accumulated transport around closed admissible
circuits.

A nontrivial closed-circuit accumulation would provide a natural
candidate for a discrete circulation or winding quantity.

That possibility remains conditional.

No theorem in this paper identifies unresolved sigma circulation with:

\begin{itemize}

\item the signed \(G_{30}\to G_{15}\) covering cocycle;

\item the twelve-section binary closure class;

\item either certified integer parent;

\item the primitive registered-history species \(r_1\);

\item or Thalean time
\[
\tau_{\mathrm T}=\operatorname{wind}_{r_1}.
\]

\end{itemize}

Those structures have different domains and independent provenance.

An explicit crosswalk would be required before any identification is
permitted.

\subsection{Compatibility failure is not saturation}

The distinction remains decisive.

A proposed global inter-face transport assignment may fail because no
collection of locally admissible seam choices satisfies all of its
closure constraints.

That is global incompatibility for that construction.

Relational saturation is different.  It concerns an already admitted
carrier or registration in which future freedom or projected
resolving power saturates while complete distinction remains.

Therefore
\[
\text{global incompatibility}
\neq
\text{relational saturation}.
\]

A failed closure search cannot be reinterpreted as evidence for
saturation.

Nor does a failed or unfinished sigma search invalidate saturation
proved on another admitted finite carrier.

\subsection{Relation to the finite saturation theorem}

Earlier versions of this paper treated a globally closed inter-face
sigma as a prerequisite for any native relational-saturation theorem.

That prerequisite is no longer required.

The finite theorem of Section 8 is carried by an independently
certified native registered carrier with exact successor,
completion, and retained-context laws.

Its saturation mechanism is
\[
8\longrightarrow2\longrightarrow1,
\]
together with retained distinction and
\[
\Delta^{\mathsf T}\Delta=\frac34H.
\]

The unresolved sigma problem therefore has a narrower status:

\begin{quote}
Sigma remains an open candidate for a stronger global inter-face
transport theory, not a missing premise of the native finite
relational-saturation theorem.
\end{quote}

\subsection{Current boundary}

At the present boundary:

\[
\text{native registered carrier}
\quad
\text{established},
\]

\[
\text{native finite relational saturation}
\quad
\text{established},
\]

\[
\text{local seam law}
\quad
\text{established},
\]

\[
E\text{-closed local vocabulary}
\quad
\text{established},
\]

\[
\text{deeper partial inter-face closure}
\quad
\text{observed},
\]

\[
\text{global inter-face sigma}
\quad
\text{open}.
\]

No physical consequence follows from the open sigma frontier.
